\section{Introducción: de Model a Agent}

Lin Junyang (林俊旸, líder del equipo de Qwen) presentó, en un evento del laboratorio conjunto de Tsinghua y Peking University, los avances técnicos de Qwen (通义千问) en 2025. El título de la ponencia cambió de ``Towards a Generalist Model'' a ``Towards a Generalist Agent''; este cambio refleja el cambio de paradigma de todo el campo de la IA: de ``entrenar modelos más fuertes'' a ``construir agentes capaces de usar herramientas de forma autónoma''.

\begin{importantbox}{¿Por qué de Model a Agent?}
Lin Junyang sostiene que Agent es un concepto más amplio que Model. Una de las diferencias centrales entre los seres humanos y los animales es \textbf{la capacidad de usar herramientas de forma autónoma}. El cambio actual en los paradigmas de entrenamiento (en particular la introducción del aprendizaje por refuerzo) le da a la IA esa misma posibilidad: en cuanto se resuelvan el razonamiento y la evaluación, el modelo podrá entrenarse para ejecutar todo tipo de tareas, ya sea como Digital Agent o como Physical Agent.
\end{importantbox}

\begin{knowledgebox}{Panorama del ecosistema de código abierto de Qwen}
\begin{itemize}
    \item El camino del código abierto comenzó el 3 de agosto de 2023
    \item Punto de acceso para probarlo: chat.qwen.ai (agrega modelos de código cerrado y abierto)
    \item Plataformas de publicación de modelos: Hugging Face y ModelScope (contenido sincronizado)
    \item GitHub se usa sobre todo para scripts; el blog técnico sustituye a los artículos (por ejemplo, la arquitectura de Qwen3-Next)
    \item La mascota es un capibara (Capybara), con un nombre de estilo desenfadado
\end{itemize}
\end{knowledgebox}

\subsection{Resumen de la sección}

La visión del equipo de Qwen es construir un ``Multimodal Foundation Agent'': un agente general capaz de ver (visión), oír (voz), hablar (generación) y actuar (uso de herramientas). Esta visión viene desde la investigación multimodal de 2020 y continúa hasta hoy.

